\section{2018级工科数学分析下 A卷}

\subsection{资料来源}
\begin{itemize}
  \item 试题：2018 级工科数学分析（二）A 卷，已导出页面图校正公式。
  \item 答案：A 卷附解答文件；抽取缺页处据题面补写解析。
\end{itemize}

\subsection{试题}

\paragraph*{试卷信息}
诚信应考，考试作弊将带来严重后果！

华南理工大学本科生期末考试《工科数学分析（二）》A卷，2018--2019学年第二学期。

注意事项：开考前请将密封线内各项信息填写清楚；所有答案请直接答在试卷上；考试形式：闭卷；本试卷共5大题，满分100分，考试时间120分钟。评阅教师请在试卷袋上评阅栏签名。

\paragraph*{一、填空题（每小题3分，共15分）}
\begin{enumerate}
  \item 微分方程 \(y''-y=2x\mathrm{e}^x\) 的特解形式为 \(y^*=x(ax+b)\mathrm{e}^x\)。
  \item 设函数 \(u(x,y,z)=\mathrm{e}^{x^2+y^2+z^2}\)，则 \(\operatorname{div}(\operatorname{grad}u)=\) \underline{\hspace{5cm}}。
  \item 设 \(L\) 是任意一条光滑的闭曲线，取逆时针方向，则
  \[
    \oint_L(2xy+\cos x)\,\mathrm{d}x+(x^2+\sin y)\,\mathrm{d}y=
  \]
  \underline{\hspace{5cm}}。
  \item 若级数 \(\displaystyle \sum_{n=1}^{\infty}\frac{(-1)^n(n+1)}{n^p}\) 条件收敛，则常数 \(p\) 的取值范围为 \underline{\hspace{5cm}}。
  \item 设 \(f(x)=\dfrac{x^2}{2}\)，\(0\le x<\pi\)，\(S(x)\) 为 \(f(x)\) 的展成以 \(2\pi\) 为周期的正弦级数的和函数，则
  \(S\left(-\dfrac{\pi}{2}\right)=\) \underline{\hspace{5cm}}。
\end{enumerate}

\paragraph*{二、计算题（每小题8分，共32分）}
\begin{enumerate}
  \item 设 \(z=f\left(x,\dfrac{y^2}{x}\right)\)，其中函数 \(f\) 有二阶连续的偏导数，求
  \(\dfrac{\partial^2 z}{\partial x\partial y}\)。

  \item 计算三重积分
  \[
    \iiint_\Omega \left(x+\frac y2-z\right)^2\,\mathrm{d}x\mathrm{d}y\mathrm{d}z,
  \]
  其中 \(\Omega\) 是曲面 \(x^2+\dfrac{y^2}{4}+z^2=1\) 所围成的区域。

  \item 求密度均匀的曲面 \(\Sigma:z=\sqrt{a^2-x^2-y^2}\) 的质心坐标。

  \item 求微分方程 \(y'\cot x=y\ln y\) 的通解。
\end{enumerate}

\paragraph*{三、解答题（每小题10分，共30分）}
\begin{enumerate}
  \item 设曲线积分
  \[
    \int_L xy^2\,\mathrm{d}x+y\varphi(x)\,\mathrm{d}y
  \]
  与积分路径无关，其中 \(\varphi(x)\) 有一阶连续导数且 \(\varphi(0)=0\)。求 \(\varphi(x)\)，并计算曲线积分
  \[
    \int_{(0,0)}^{(1,1)}xy^2\,\mathrm{d}x+y\varphi(x)\,\mathrm{d}y.
  \]

  \item 计算
  \[
    I=\oiint_\Sigma
    \frac{x\,\mathrm{d}y\mathrm{d}z+y\,\mathrm{d}z\mathrm{d}x+z\,\mathrm{d}x\mathrm{d}y}
    {(a^2x^2+b^2y^2+c^2z^2)^{3/2}},
  \]
  其中 \(a,b,c>0\)，\(\Sigma:x^2+y^2+z^2=1\)，取外侧。

  \item 求幂级数 \(\displaystyle \sum_{n=1}^{\infty}\frac{2n-1}{n!}x^{2n}\) 的收敛域及和函数。
\end{enumerate}

\paragraph*{四、证明题（每小题7分，共14分）}
\begin{enumerate}
  \item 设 \(f(x)\in C[0,a]\)，证明：
  \[
    2\int_0^a f(x)\,\mathrm{d}x\int_x^a f(y)\,\mathrm{d}y
    =\left[\int_0^a f(x)\,\mathrm{d}x\right]^2.
  \]
  \item 证明函数项级数
  \[
    \sum_{n=1}^{\infty}\frac{\arctan(nx)}{n^2}
  \]
  在 \((-\infty,+\infty)\) 上一致收敛。
\end{enumerate}

\paragraph*{五、应用题（本题9分）}
在椭圆 \(x^2+4y^2=4\) 的第一象限部分上求一点，使其到直线 \(2x+3y-6=0\) 的距离最短。

\subsection{答案与解析}

\paragraph*{一、填空题}
\begin{enumerate}
  \item 由于 \(1\) 是特征方程 \(r^2-1=0\) 的单根，故特解形式为
  \[
    y^*=x(ax+b)\mathrm{e}^x.
  \]
  \item
  \[
    \operatorname{div}(\operatorname{grad}u)=\Delta u
    =\left[6+4(x^2+y^2+z^2)\right]\mathrm{e}^{x^2+y^2+z^2}.
  \]
  \item 由 Green 公式，\(Q_x-P_y=2x-2x=0\)，故积分为 \(0\)。
  \item 交错收敛要求 \((n+1)/n^p\to0\)，即 \(p>1\)；绝对收敛要求 \(\sum n^{1-p}\) 收敛，即 \(p>2\)。故条件收敛范围为
  \[
    1<p\le2.
  \]
  \item 正弦级数对应奇延拓，故
  \[
    S\left(-\frac{\pi}{2}\right)=-f\left(\frac{\pi}{2}\right)=-\frac{\pi^2}{8}.
  \]
\end{enumerate}

\paragraph*{二、计算题}
\begin{enumerate}
  \item 记 \(u=x,\ v=y^2/x\)，则 \(z=f(u,v)\)。先求
  \[
    z_y=f_2\frac{2y}{x}.
  \]
  再对 \(x\) 求导，得到
  \[
    \frac{\partial^2z}{\partial x\partial y}
    =\frac{2y}{x}f_{12}-\frac{2y^3}{x^3}f_{22}-\frac{2y}{x^2}f_2,
  \]
  其中 \(f_2,f_{12},f_{22}\) 均在 \(\left(x,\dfrac{y^2}{x}\right)\) 处取值。

  \item 展开平方后，由椭球关于三个坐标面对称，交叉项积分均为 \(0\)，所以原积分等于
  \[
    \iiint_\Omega\left(x^2+\frac{y^2}{4}+z^2\right)\,\mathrm{d}V.
  \]
  作变换
  \[
    x=r\sin\varphi\cos\theta,\quad
    y=2r\sin\varphi\sin\theta,\quad
    z=r\cos\varphi,\quad J=2r^2\sin\varphi.
  \]
  因此
  \[
    I=\int_0^{2\pi}\!\!\int_0^\pi\!\!\int_0^1 r^2\cdot2r^2\sin\varphi\,\mathrm{d}r\,\mathrm{d}\varphi\,\mathrm{d}\theta
    =\frac{8\pi}{5}.
  \]

  \item 由对称性，\(\bar x=\bar y=0\)。上半球面面积为 \(2\pi a^2\)，且
  \[
    \iint_\Sigma z\,\mathrm{d}S
    =\int_0^{2\pi}\!\!\int_0^{\pi/2}a\cos\varphi\cdot a^2\sin\varphi\,\mathrm{d}\varphi\,\mathrm{d}\theta
    =\pi a^3.
  \]
  故质心为
  \[
    \left(0,0,\frac a2\right).
  \]

  \item 当 \(y>0\) 时，方程化为
  \[
    \frac{y'}{y\ln y}=\tan x.
  \]
  令 \(u=\ln y\)，则
  \[
    \frac{u'}u=\tan x,\qquad \ln|u|=-\ln|\cos x|+C.
  \]
  因此通解为
  \[
    \ln y=C\sec x,\qquad y=\mathrm{e}^{C\sec x}.
  \]
  其中 \(C=0\) 给出常值解 \(y=1\)。
\end{enumerate}

\paragraph*{三、解答题}
\begin{enumerate}
  \item 令 \(P=xy^2\)，\(Q=y\varphi(x)\)。路径无关要求
  \[
    P_y=Q_x,\qquad 2xy=y\varphi'(x),
  \]
  故 \(\varphi'(x)=2x\)。由 \(\varphi(0)=0\)，得
  \[
    \varphi(x)=x^2.
  \]
  势函数可取
  \[
    F(x,y)=\frac12x^2y^2,
  \]
  所以
  \[
    \int_{(0,0)}^{(1,1)}xy^2\,\mathrm{d}x+yx^2\,\mathrm{d}y
    =F(1,1)-F(0,0)=\frac12.
  \]

  \item 在单位球面上，外法向量为 \((x,y,z)\)，故
  \[
    x\,\mathrm{d}y\mathrm{d}z+y\,\mathrm{d}z\mathrm{d}x+z\,\mathrm{d}x\mathrm{d}y=\mathrm{d}S.
  \]
  于是
  \[
    I=\iint_{x^2+y^2+z^2=1}
    \frac{\mathrm{d}S}{(a^2x^2+b^2y^2+c^2z^2)^{3/2}}.
  \]
  由线性变换 \(X=ax,\ Y=by,\ Z=cz\) 对单位球面面积元的标准公式可得
  \[
    I=\frac{4\pi}{abc}.
  \]

  \item 令 \(t=x^2\)。因为
  \[
    \sum_{n=1}^{\infty}\frac{nt^n}{n!}=t\mathrm{e}^t,\qquad
    \sum_{n=1}^{\infty}\frac{t^n}{n!}=\mathrm{e}^t-1,
  \]
  所以
  \[
    \sum_{n=1}^{\infty}\frac{2n-1}{n!}x^{2n}
    =2x^2\mathrm{e}^{x^2}-(\mathrm{e}^{x^2}-1)
    =(2x^2-1)\mathrm{e}^{x^2}+1.
  \]
  收敛域为 \((-\infty,+\infty)\)。
\end{enumerate}

\paragraph*{四、证明题}
\begin{enumerate}
  \item 令
  \[
    F(x)=\int_x^a f(y)\,\mathrm{d}y.
  \]
  则 \(F'(x)=-f(x)\)，从而
  \[
    2\int_0^a f(x)\,\mathrm{d}x\int_x^a f(y)\,\mathrm{d}y
    =2\int_0^a f(x)F(x)\,\mathrm{d}x
    =-2\int_0^aF(x)F'(x)\,\mathrm{d}x.
  \]
  因此
  \[
    -2\int_0^aF(x)F'(x)\,\mathrm{d}x
    =F(0)^2-F(a)^2
    =\left[\int_0^a f(x)\,\mathrm{d}x\right]^2.
  \]

  \item 对任意 \(x\in(-\infty,+\infty)\)，
  \[
    \left|\frac{\arctan(nx)}{n^2}\right|\le \frac{\pi}{2n^2}.
  \]
  由于 \(\sum_{n=1}^{\infty}\dfrac{\pi}{2n^2}\) 收敛，由 Weierstrass 判别法，原函数项级数在 \((-\infty,+\infty)\) 上一致收敛。
\end{enumerate}

\paragraph*{五、应用题}
第一象限椭圆上有 \(2x+3y\le6\)，故到直线 \(2x+3y-6=0\) 的距离最短等价于使 \(2x+3y\) 最大。设
\[
  F(x,y)=2x+3y,\qquad x^2+4y^2=4.
\]
由 Lagrange 乘子法，
\[
  2=2\lambda x,\qquad 3=8\lambda y.
\]
解得
\[
  x=\frac85,\qquad y=\frac35.
\]
因此所求点为
\[
  \left(\frac85,\frac35\right).
\]
